\documentclass[10pt,a4paper]{article}

\usepackage[spanish,es-nodecimaldot]{babel}
\usepackage[T1]{fontenc}
\usepackage[utf8]{inputenc}
\usepackage{lmodern}
\usepackage[margin=1.35cm]{geometry}
\usepackage{xcolor}
\usepackage{enumitem}
\usepackage{listings}
\usepackage{hyperref}
\usepackage{multicol}
\usepackage{microtype}
\usepackage{fancyhdr}

\definecolor{lpmdblue}{HTML}{174A73}
\definecolor{lpmdlight}{HTML}{EEF5FA}
\definecolor{lpmdgreen}{HTML}{24735C}
\definecolor{codegray}{HTML}{F3F4F5}

\hypersetup{colorlinks=true,urlcolor=lpmdblue,linkcolor=lpmdblue}
\setlength{\parindent}{0pt}
\setlength{\parskip}{5pt}
\setlength{\columnsep}{0.85cm}
\setlist[itemize]{leftmargin=*,itemsep=2pt,topsep=3pt}
\pagestyle{fancy}
\fancyhf{}
\renewcommand{\headrulewidth}{0pt}
\fancyfoot[C]{\footnotesize LPMD 2 --- Guía rápida de instalación y primera ejecución}

\lstdefinestyle{terminal}{
  basicstyle=\ttfamily\footnotesize,
  backgroundcolor=\color{codegray},
  frame=single,
  rulecolor=\color{lpmdblue!35},
  framerule=0.4pt,
  breaklines=true,
  columns=fullflexible,
  keepspaces=true,
  showstringspaces=false,
  aboveskip=6pt,
  belowskip=7pt
}

\newcommand{\sectiontitle}[1]{%
  \vspace{5pt}{\large\bfseries\color{lpmdblue}#1}\par
  \vspace{2pt}\color{lpmdblue!35}\hrule\color{black}\vspace{6pt}
}

\newcommand{\notice}[1]{%
  \fcolorbox{lpmdgreen!55}{lpmdgreen!7}{%
    \begin{minipage}{0.94\linewidth}\small #1\end{minipage}}}

\begin{document}

\begin{center}
  {\Huge\bfseries\color{lpmdblue} Instalar y probar LPMD 2}\par
  \vspace{2pt}
  {\large Guía rápida para Linux}\par
  \vspace{5pt}
  \color{lpmdblue!35}\rule{\linewidth}{1.2pt}
\end{center}

\begin{multicols}{2}

\sectiontitle{0. Obtener el código}

Clona el repositorio y entra en su directorio. Los pasos siguientes deben
ejecutarse desde allí:

\begin{lstlisting}[style=terminal]
git clone https://github.com/jperaltac/lpmd2.git
cd lpmd2
\end{lstlisting}

\sectiontitle{1. Preparar el sistema}

Instala las dependencias de compilación:

\begin{lstlisting}[style=terminal]
./scripts/install_requirements.sh
\end{lstlisting}

Se requiere CMake 3.16 o posterior, un compilador C++17 y los archivos de
desarrollo de zlib. OpenGL/GLUT es opcional para los plugins gráficos.

\sectiontitle{2. Compilar}

Configura una compilación optimizada y construye todos los objetivos:

\begin{lstlisting}[style=terminal]
cmake -S . -B build \
  -DCMAKE_BUILD_TYPE=Release
cmake --build build -j
\end{lstlisting}

Los ejecutables quedan temporalmente en \texttt{build/bin} y los plugins en
\texttt{build/lib}.

\sectiontitle{3. Instalar}

\textbf{Opción A: instalación para todo el sistema}

El prefijo predeterminado es \texttt{/usr/local}; normalmente requiere
privilegios administrativos:

\begin{lstlisting}[style=terminal]
sudo cmake --install build
\end{lstlisting}

\textbf{Opción B: instalación solo para el usuario}

No requiere \texttt{root}. El prefijo debe definirse al configurar:

\begin{lstlisting}[style=terminal]
cmake -S . -B build-user \
  -DCMAKE_BUILD_TYPE=Release \
  -DCMAKE_INSTALL_PREFIX="$HOME/.local"
cmake --build build-user -j
cmake --install build-user
\end{lstlisting}

Si es necesario, habilita los ejecutables instalados en la sesión actual:

\begin{lstlisting}[style=terminal]
export PATH="$HOME/.local/bin:$PATH"
\end{lstlisting}

\columnbreak

\sectiontitle{4. Comprobar la instalación}

Comprueba que el sistema encuentre el programa:

\begin{lstlisting}[style=terminal]
command -v lpmd
lpmd --help
\end{lstlisting}

La instalación incluye:

\begin{itemize}
  \item los ejecutables \texttt{lpmd}, \texttt{lpmd-analyzer},\newline
        \texttt{lpmd-converter} y \texttt{lpmd-visualizer};
  \item la biblioteca \texttt{liblpmd} y sus cabeceras C++;
  \item los módulos cargables en \texttt{lib/lpmd};
  \item integración con CMake y \texttt{pkg-config}.
\end{itemize}

\sectiontitle{5. Ejecutar un ejemplo simple}

El ejemplo \texttt{smoke.control} crea cuatro átomos de argón, ejecuta un paso
sin fuerzas y escribe \texttt{smoke.xyz}. Para mantener limpio el repositorio:

\begin{lstlisting}[style=terminal]
mkdir -p /tmp/lpmd-smoke
cd /tmp/lpmd-smoke
lpmd /ruta/al/repositorio/lpmd/examples/\
smoke.control
\end{lstlisting}

Reemplaza \texttt{/ruta/al/repositorio} por la ruta real del checkout. Una
ejecución correcta muestra:

\begin{lstlisting}[style=terminal]
SIMULATION FINISHED
\end{lstlisting}

y crea \texttt{/tmp/lpmd-smoke/smoke.xyz}. Puedes inspeccionar sus primeras
líneas mediante:

\begin{lstlisting}[style=terminal]
head smoke.xyz
\end{lstlisting}

\sectiontitle{Solución rápida de problemas}

\begin{itemize}
  \item Si \texttt{lpmd} no aparece, revisa \texttt{PATH} o ejecútalo con su
        ruta completa, por ejemplo \texttt{\$HOME/.local/bin/lpmd}.
  \item Si usas plugins externos, agrega sus directorios a
        \texttt{LPMD\_PATH}, separados por dos puntos.
  \item Para cambiar el prefijo, vuelve a ejecutar la configuración de CMake;
        no basta con cambiarlo solamente durante \texttt{cmake --install}.
\end{itemize}

\vfill
\notice{\textbf{Consejo:} para actualizar la instalación después de modificar
el código, repite \texttt{cmake --build build} y luego
\texttt{sudo cmake --install build}.}

\end{multicols}

\end{document}
